\chapter{The Internals of Mathilda}
\label{ch:internals}

Everything in this book so far has looked at Mathilda from the outside: you type an
expression, and an answer comes back. This chapter turns the machine over and looks
at the works. It is the one place in the book where we get into the C-level detail
of how the system is actually built---how an expression is laid out in memory, what
the evaluator does on every step and in exactly what order, how the pattern matcher
searches, how memory is owned and reclaimed, and how the parser, printer and the
interactive loop bracket the whole thing.

Appendix~\ref{app:architecture} is the aerial photograph: the pipeline in one
figure, each subsystem named in a paragraph. This chapter is the walk on the
ground. If you have not read that appendix, Figure~\ref{fig:pipeline} there is worth
a glance first---it shows how the five components fit together and it is the frame we
now fill in. Where the appendix says ``the evaluator reads the head's attributes,''
this chapter tells you which field it reads, when, and what it costs.

Two conventions for what follows. The \emph{transcripts} are, as everywhere in this
book, real output from the running binary---the internals are observable from the
REPL more often than you might expect, and we lean on that. The \emph{C excerpts},
tinted blue and titled with their source file, are the exception: they are quoted by
hand from the source tree to show a data structure or a control path, abridged for
the page and \emph{not} run by the build the way the transcripts are. Treat them as a
faithful sketch and open the named file for the whole truth.

\section{Expression trees}
\label{sec:int-expr}

The foundational decision, the one from which everything else follows, is that
\emph{every value is the same kind of object}: an \texttt{Expr}. A number, a symbol,
a string, a list, a rule, a pattern, a matrix, a held piece of code---all of them are
\texttt{Expr} nodes, and the whole system is machinery for building, comparing,
rewriting and freeing them. You saw this from the user's side in
Chapter~\ref{ch:datastructures}; here is the node itself.

An \texttt{Expr} is a \emph{tagged union}. A single enum says which kind of value the
node holds, and a union holds the payload for that kind:

\begin{ccode}[src/expr.h (abridged)]
typedef enum {
    EXPR_INTEGER, EXPR_REAL, EXPR_SYMBOL, EXPR_STRING,
    EXPR_FUNCTION, EXPR_BIGINT, EXPR_NDARRAY, EXPR_COMPILED
#ifdef USE_MPFR
    , EXPR_MPFR
#endif
} ExprType;

typedef struct Expr {
    ExprType type;
    unsigned refcount;            /* shared ownership / copy-on-write         */
    uint64_t last_evaluated_at;   /* eval-clock stamp; top bit = GROUND flag  */
    uint64_t hash_cache;          /* memoized structural hash; 0 = unset      */
    union {
        int64_t integer;
        double  real;
        struct { char* name; struct SymbolDef* def; } symbol;
        char*   string;
        struct {
            struct Expr*  head;
            struct Expr** args;
            size_t        arg_count;
            struct AssocIndex*  index;         /* Association key->position   */
            struct SymSetCache* symset_cache;  /* symbol set of the subtree  */
        } function;
        mpz_t bigint;                          /* GMP arbitrary-precision int */
        NDArrayData ndarray;                   /* dense machine-number buffer */
        struct CompiledFunction* compiled;
#ifdef USE_MPFR
        mpfr_t mpfr;                           /* arbitrary-precision real    */
#endif
    } data;
} Expr;
\end{ccode}

The classic six kinds are the ones you meet first---integer, real, symbol, string,
function, and the arbitrary-precision \texttt{EXPR\_BIGINT}. Two more,
\texttt{EXPR\_NDARRAY} and \texttt{EXPR\_COMPILED}, are the dense numeric buffer of
\S\ref{sec:int-packed} and the compiled-code object of
Chapter~\ref{ch:compilation}; a ninth, \texttt{EXPR\_MPFR}, appears only in builds
with extended-precision reals. All of the last three are \emph{atomic} to the
evaluator---it never looks inside them.

A \emph{function} node is the compound case, and it is the interesting one: a head
(itself an \texttt{Expr*}) and a heap array of argument pointers with an explicit
count. There is no separate capacity field---the array is allocated to exactly
\texttt{arg\_count} entries---which matters later, when we pool these arrays by
length (\S\ref{sec:int-memory}). The compound form \(\mathtt{f[x, y]}\) is a function
node whose head is the symbol \texttt{f}; a list is a function node with head
\texttt{List}. You can X-ray any surface syntax into this tree with \B{FullForm}, and
confirm each node's kind with \B{Head} and \B{AtomQ}:

\mtranscript{10-internals/expr-fullform}

Read those as the tree behind the syntax. A sum is \texttt{Plus[a, Times[b, c]]}; a
fraction is a single \texttt{Rational} node; a division is a \texttt{Times} against a
\texttt{Power} with exponent \(-1\); a complex literal is \texttt{Complex[2, 3]}; and
even a pattern, \(\mathtt{x\_}\), is an ordinary expression,
\texttt{Pattern[x, Blank[]]}. \B{Head} of an integer is \texttt{Integer}, of a list
is \texttt{List}: atoms carry a head too, which is why the same tools---\B{Part},
\B{Map}, \B{Replace}---work on everything. \B{AtomQ} draws the one structural line
that matters internally: an atom has no arguments to descend into, a function node
does.

\subsection*{The three metadata words}

Before the union sit three machine words that are not part of the value at all. They
are pure acceleration, invisible to equality and ordering, and understanding them
explains most of what makes Mathilda fast.

\texttt{refcount} is the reference count: shared ownership, discussed in full in
\S\ref{sec:int-memory}. \texttt{last\_evaluated\_at} is a timestamp the evaluator
stamps on a node when it reaches a fixed point, so a later evaluation of the same
node can return instantly (\S\ref{sec:int-eval}); its top bit is a flag for
loop-invariant values. \texttt{hash\_cache} memoizes the node's structural hash the
first time anything asks for it, turning the repeated subtree walks in \texttt{Plus}
grouping, \texttt{Association} indexing and set operations into
\(O(1)\) reads.\calloutnote{Under the hood}{Because these three fields are metadata,
not value, \texttt{expr\_eq}, \texttt{expr\_hash} and \texttt{expr\_compare} all
ignore them---two nodes are equal when their \emph{structure} matches, whatever their
refcounts or stamps. The flip side is a strict rule the whole tree obeys: a cached
hash is only safe because a shared node is never mutated in place (\S\ref{sec:int-memory}),
and any code that does rewrite a node's fields must reset \texttt{hash\_cache} to 0
first. The \texttt{MATHILDA\_HASH\_VERIFY} build recomputes and asserts every cache
hit to catch a missed reset.}

Integers show a fourth theme---representations that adapt. A value begins life as a
machine \texttt{int64}, and is promoted to a GMP bignum only when an operation would
overflow, demoting back when a result fits again. There is only ever ``an integer,''
never a user-visible big-integer mode:

\mtranscript{10-internals/expr-bignum}

\(2^{62}\) is a machine word; \(2^{63}\) has already overflowed into a
\texttt{EXPR\_BIGINT}; \(100!\) is a 158-digit bignum---and \B{Head} says
\texttt{Integer} for all of them. The promotion is a change of internal
representation, not of type, and it is completely invisible from the language.

\section{The evaluator}
\label{sec:int-eval}

The evaluator is the heart of the system, and its governing idea is
\emph{fixed-point rewriting}: \texttt{evaluate(e)} calls \texttt{evaluate\_step(e)}
over and over, each call rewriting the expression by one level, and stops when a step
changes nothing. There is no separate ``simplifier'' pass and no explicit chaining of
rules; composition is a consequence of running the step to a fixed point. You can
watch the process directly with \B{Trace}, which records every intermediate expression:

\mtranscript{10-internals/eval-trace}

A definition and an arithmetic fold each show the shape of it---\texttt{f[3]} becomes
\texttt{3\^{}2} becomes \texttt{9}; the three-term sum collapses to \texttt{9} in one
step. And because a step just rewrites and the loop merely repeats, composition needs
no plumbing at all: define one function in terms of another and the answer falls out
over successive passes.

\mtranscript{10-internals/eval-fixedpoint}

Each element of the \B{Trace} list is one intermediate expression on the way to the
answer: \texttt{g[10]} rewrites to \texttt{h[10]}, then to \texttt{10 + 1}, then to
\texttt{11}, at which point a step produces nothing new and the loop halts. No code in
\texttt{g} or \texttt{h} knows about the other; the loop is the only glue.

\subsection*{One step, in order}

The real content is what \texttt{evaluate\_step} does for a function call
\(\mathtt{f[a_1, \dots, a_n]}\). The appendix gives the eight-line summary; the true
sequence, in \texttt{src/eval.c}, is longer and every stage earns its place.
Figure~\ref{fig:evalstep} lists it in order.

\begin{figure}[h]
\begin{framed}
\small
\begin{verbatim}
   evaluate_step( f[a1, ..., an] ):

    1.  evaluate the head f
    2.  resolve f's SymbolDef once, cache it on the head node;
        read attributes as (base_attributes | attributes)  -- O(1)
    3.  evaluate each argument UNLESS a Hold attribute holds that slot
          (Evaluate[] punches through a hold; atoms skip the recursion)
    4.  reassemble f[...] from the evaluated arguments
    5.  flatten Sequence[...] arguments        (unless SequenceHold)
    6.  strip Nothing from a List
    7.  strip Unevaluated[...] wrappers in non-held slots
   7.5  transparency gate: materialize any packed buffer whose head
          is not packed-aware                  (see 10.9)
    8.  Flat      -- flatten nested f[f[...]]   if ATTR_FLAT
    9.  Listable  -- thread over list arguments if ATTR_LISTABLE
   10.  Orderless -- sort arguments canonically if ATTR_ORDERLESS
   11.  try f's DOWNVALUES  (user/library rules)  -- before the builtin
   12.  try f's NDArray element-wise machine kernel, if registered
   13.  call f's builtin C function, if any
   14.  Set / SetDelayed and the composite-head cases (pure Function, ...)
\end{verbatim}
\end{framed}
\caption{One rewrite level. Steps 8--10 apply the associative, threading and
commutative attributes in that fixed order; the surprise for a newcomer is step 11,
where user \texttt{DownValues} are tried \emph{before} the builtin C code, so a rule
you write can override a primitive.}
\label{fig:evalstep}
\end{figure}

Three things about this sequence are worth dwelling on.

\emph{Attributes are read once, and cheaply.} Step~2 resolves the head symbol to its
definition record and caches the pointer on the head node, so subsequent touches skip
the symbol-table lookup entirely. The attributes themselves are a single field read
and a bitwise OR---\texttt{base\_attributes} (the immutable floor a builtin is born
with) combined with \texttt{attributes} (whatever the user has since set). This
replaced an older design that walked a 140-entry table of name comparisons on every
call. The attribute bits are the small vocabulary that steers the generic loop; the
full set is in Table~\ref{tab:attrs}.

\emph{Holding happens in step~3, bit by bit.} Whether an argument is evaluated is
decided per slot, by three lines:

\begin{ccode}[src/eval.c (the argument loop)]
bool hold = hold_all_complete;                        /* HoldAllComplete: every slot */
if (i == 0 && (attrs & ATTR_HOLDFIRST)) hold = true;  /* HoldFirst: slot 0           */
if (i >  0 && (attrs & ATTR_HOLDREST))  hold = true;  /* HoldRest:  slots 1..n       */
\end{ccode}

\texttt{HoldAll} needs no case of its own: it is simply \texttt{HoldFirst \& HoldRest}
as a combined mask. Everything about non-standard evaluation---\texttt{Table} not
evaluating its iterator, \texttt{If} not evaluating the branch it will not take,
\texttt{Set} not evaluating its left side---reduces to which of these bits the head
carries. That is the whole subject of \S\ref{sec:int-nonstd}.

\emph{DownValues come before the builtin.} Step~11 is ahead of step~13 on purpose: a
rule the user or a library has installed on \texttt{f} is tried before \texttt{f}'s
own C code, so mathematics expressed as a rewrite rule (the integral tables, the
derivative rules of \S\ref{sec:int-boot}) can extend or override a primitive without
touching C.

The attribute bits do all the steering, and you can read a head's bits from the
language with \B{Attributes} and change them with \B{SetAttributes}:

\mtranscript{10-internals/eval-attributes}

\texttt{Plus} carries five attributes at once; \texttt{Sin} only two; \texttt{Table}
is \texttt{HoldAll}; \texttt{Set} is \texttt{HoldFirst}. Watch the effect of each: an
\texttt{Orderless} head sorts \texttt{b + a + c} into canonical order; adding
\texttt{Flat} to \texttt{f} makes \texttt{f[f[a, b], c]} flatten to
\texttt{f[a, b, c]}; adding \texttt{Listable} to \texttt{h} makes it thread over a
list. Nothing about the main loop changed between those heads---only their bits did.

\begin{table}[h]
\centering
\small
\begin{tabular}{lll}
\hline
Attribute & Bit & Effect \\
\hline
\texttt{HoldFirst} & \texttt{1<<0} & do not evaluate argument 0 \\
\texttt{HoldRest} & \texttt{1<<1} & do not evaluate arguments 1..n \\
\texttt{HoldAll} & (both) & do not evaluate any argument \\
\texttt{HoldAllComplete} & \texttt{1<<2} & hold all, and suppress Sequence/Flat/Unevaluated \\
\texttt{Flat} & \texttt{1<<3} & associative: flatten nested same-head calls \\
\texttt{Orderless} & \texttt{1<<4} & commutative: sort arguments canonically \\
\texttt{Listable} & \texttt{1<<5} & thread element-wise over list arguments \\
\texttt{NumericFunction} & \texttt{1<<6} & numeric under \texttt{N}; threads over intervals \\
\texttt{Protected} & \texttt{1<<7} & values and attributes cannot be changed \\
\texttt{OneIdentity} & \texttt{1<<8} & \(\mathtt{f[x]}\equiv\mathtt{x}\) for matching (\S\ref{sec:int-match}) \\
\texttt{SequenceHold} & \texttt{1<<13} & do not flatten \texttt{Sequence} in arguments \\
\texttt{Constant} & \texttt{1<<14} & derivative is zero \\
\hline
\end{tabular}
\caption{The core attribute bits, from \texttt{src/attr.h}. \texttt{Locked},
\texttt{Temporary} and the \texttt{NHold*} family occupy the remaining bits; bit 11,
once \texttt{ReadProtected}, was removed---Mathilda never hides a definition.}
\label{tab:attrs}
\end{table}

\subsection*{Knowing when to stop}

A fixed-point loop needs to know when to quit, and doing that naively---comparing the
whole tree before and after every step---would be quadratic on deep expressions.
Mathilda avoids it with a cheap boolean: \texttt{evaluate\_step} reports whether it
actually rewrote anything, and only when it claims a change does the loop fall back to
a structural comparison. When a step fires no rule, the loop stops without walking the
tree at all.

Stability across \emph{calls} is handled by the \texttt{last\_evaluated\_at} stamp. A
global evaluation clock ticks forward on every definition change---every \texttt{Set},
\texttt{SetDelayed}, \texttt{Clear} or \texttt{SetAttributes}. When
\texttt{evaluate} reaches a fixed point it stamps the result with the current clock;
a later \texttt{evaluate} of a node already carrying the current stamp returns an
inc-ref'd view immediately, skipping the loop and all of \texttt{evaluate\_step}. A
single definition change, by bumping the clock, invalidates every cached evaluation in
the system at once.\calloutnote{Under the hood}{The stamp's top bit is a second,
finer cache. A node built entirely from literals under one of six pure structural
constructors (\texttt{List}, \texttt{Association}, \texttt{Rule},
\texttt{RuleDelayed}, \texttt{Complex}, \texttt{Rational}) is marked \emph{ground}:
its value cannot depend on any symbol's definition, so it survives clock churn and can
be re-validated as a fixed point in \(O(1)\) even after an ordinary assignment. This
is what keeps a loop that rebuilds a constant list on every iteration from
re-evaluating it every time.}

Two limits bound the process, and they are deliberately asymmetric.
\(\mathtt{\$RecursionLimit}\) caps how deep evaluation may nest; it is a real,
user-settable symbol (default 1024, with a floor of 20), and on overflow the offending
expression is returned wrapped in \texttt{Hold} so it cannot immediately re-enter.
\(\mathtt{\$IterationLimit}\), by contrast, is \emph{not} a settable symbol at all: the
number of fixed-point iterations for a single expression is capped by a compile-time
constant (4096). It is an honest simplification---the constant is generous enough that
only a pathological non-terminating rewrite reaches it---but it is worth knowing that,
unlike in some systems, you cannot raise it from the language.

\section{Non-standard evaluation}
\label{sec:int-nonstd}

``Non-standard evaluation'' is the umbrella term for everything that departs from the
default ``evaluate all the arguments, then call the head.'' In many languages this
would need special forms baked into the interpreter. In Mathilda it is one mechanism:
the \texttt{Hold} attributes of \S\ref{sec:int-eval}, step~3. A head that wants to
control its arguments simply carries \texttt{HoldFirst}, \texttt{HoldRest} or
\texttt{HoldAll}, and the generic loop hands it the arguments unevaluated. Everything
in this section is built from that one bit-test.

\usagebox{Hold}

The primitives that manage held code are small and composable:

\mtranscript{10-internals/nonstandard-hold}

\B{Hold} keeps its contents literally---it is nothing but a head carrying
\texttt{HoldAll}, so \texttt{Hold[1 + 1]} never adds. \B{Evaluate} punches a hole in a
hold, forcing just its argument while its siblings stay held. \B{Unevaluated} is the
mirror image: it passes an argument through unevaluated and then dissolves, so
\texttt{Length[Unevaluated[1 + 2 + 3]]} sees the three-element \texttt{Plus} and
answers \texttt{3}. \B{ReleaseHold} strips a wrapper so its contents can finally run.
The distinction between \B{Set} and \B{SetDelayed} is the same mechanism: \texttt{=}
is \texttt{HoldFirst} (it holds the left side but evaluates the right side now), while
\texttt{:=} is \texttt{HoldAll} (it holds both and re-evaluates the right side on every
use). And \B{If}, carrying \texttt{HoldRest}, never touches the branch it does not
take---which is why the untaken \texttt{boom = 1/0} above caused no error and left
\texttt{boom} unassigned.

\subsection*{Scoping without a stack}

The scoping constructs are the richest users of held arguments, and they achieve their
effects in three distinct ways---worth contrasting, because the differences are
observable:

\mtranscript{10-internals/nonstandard-scoping}

\B{Module} implements lexical locals by \emph{alpha-renaming}: each local \texttt{x} is
renamed to a globally fresh symbol \texttt{x\$}\textit{n}, and you can see the counter
advance---the two identical \texttt{Module} calls return \texttt{Hold[x\$1]} and
\texttt{Hold[x\$2]}. \B{Block}, by contrast, is \emph{dynamic} scoping: it does not
rename anything but temporarily saves the global value of \texttt{n}, runs the body,
and restores it---so inside the block \texttt{n} is 2 and after it \texttt{n} is 10
again. \B{With} does neither: it substitutes its constants directly into the body
before evaluation, a purely lexical rewrite. Pure functions
(\(\mathtt{\#^2\, \&}\)) bind their \texttt{Slot} arguments by the same
substitution.\calloutnote{Under the hood}{None of these three touches a call stack of
activation records---Mathilda has no such thing. \texttt{Module} works by renaming and
installing a temporary \texttt{OwnValue}; \texttt{Block} by saving and restoring a
symbol's value cell; \texttt{With} and pure functions by lexical substitution into a
held body. \texttt{Table} and the other iterators are dynamic scoping too, saving the
iterator symbol's value around the loop. The uniform ``everything is an expression''
model means a binding is just a value on a symbol, and scoping is just deciding when to
put it there and when to take it away.}

\section{The pattern matcher}
\label{sec:int-match}

Pattern matching is how Mathilda decides whether a rule applies, and it is the most
combinatorially delicate subsystem in the tree. Chapter~\ref{ch:programming} taught
the pattern \emph{language}---blanks, sequences, tests, conditions; this section is the
engine beneath it, \texttt{src/match.c}. From the language side you drive it with
\B{MatchQ} and its relatives:

\mtranscript{10-internals/matcher-basics}

Everything in that session is the matcher at work: a typed blank checks a head; a
repeated name enforces equality; sequence blanks match runs of arguments; an
\texttt{Orderless} head matches against reordered arguments; a test or condition
gates a match; an \texttt{Optional} supplies a default. We now take the mechanism
apart.

\subsection*{Bindings and backtracking}

Pattern variables are held in a \texttt{MatchEnv}: two parallel arrays, one of names
and one of bound values, with a linear scan for lookup. A handful of variables is the
common case, so a scan beats the overhead of a hash.

\begin{ccode}[src/match.h and match.c]
typedef struct MatchEnv {
    char** symbols;     /* borrowed interned names -- NOT strdup'd     */
    Expr** values;      /* the env owns a copy of each bound value     */
    size_t count;
    size_t capacity;
    bool (*callback)(struct MatchEnv*, void*);  /* ReplaceList: see every match */
    void* callback_data;
} MatchEnv;

/* Backtracking is a single integer. Bindings are only ever appended, so
 * truncating the count back to a checkpoint is a complete, cheap undo. */
static void env_rollback(MatchEnv* env, size_t saved_count) {
    while (env->count > saved_count)
        expr_free(env->values[--env->count]);
}
\end{ccode}

That \texttt{env\_rollback} is the whole backtracking discipline. At every branch
point the matcher records \texttt{env->count}, tries an alternative, and if it fails
truncates the bindings back to the saved count and tries the next. Because a name is
never overwritten---only appended after a check---rolling the count back is an exact
undo. \emph{Non-linear} patterns fall straight out of this: when the matcher is about
to bind a name that already has a value, it does not re-bind; it compares the new
candidate to the existing binding with \texttt{expr\_eq} and fails if they differ.
That single rule is why \texttt{f[x\_, x\_]} matches \texttt{f[a, a]} but not
\texttt{f[a, b]}.

The recursion itself threads a continuation---a \texttt{ParentMatch} chain
representing ``the sibling patterns still to match''---so that a deep sub-match, on
success, can carry on matching the arguments that follow it. There is no explicit work
stack; the C call stack, guarded by the same \(\mathtt{\$RecursionLimit}\) the
evaluator uses, is the stack.

\subsection*{Sequences: the combinatorial heart}

Ordinary structural matching is easy: same head, same arity, recurse pairwise. The
hard part is sequence patterns, \(\mathtt{x\_\_}\) and \(\mathtt{x\_\_\_}\), which
match \emph{runs} of arguments of unknown length. The matcher searches the partitions:
for the first sequence pattern it tries letting it consume \(k\) arguments for each
legal \(k\), and for each choice recurses on the rest. Figure~\ref{fig:backtrack}
shows the search for \texttt{f[x\_\_, y\_\_]} against \texttt{f[1, 2, 3]}.

\begin{figure}[h]
\begin{framed}
\small
\begin{verbatim}
   match  f[x__, y__]   against   f[1, 2, 3]        (each __ needs >= 1 arg)

     x__ = {1}          x__ = {1,2}         x__ = {1,2,3}
       |                   |                    |
     y__ = {2,3}  OK     y__ = {3}   OK       y__ = {}   FAIL
       ^                                        (y__ requires >= 1)
       first success returned
\end{verbatim}
\end{framed}
\caption{Sequence matching is a partition search with backtracking. Each split is
tried in turn; the first that lets every remaining pattern match wins.}
\label{fig:backtrack}
\end{figure}

Two forces make this affordable in practice. The first is a hard short-circuit: when a
sequence pattern is the \emph{last} one in the argument list, it must consume all the
remaining arguments, so there is nothing to search---\(k\) is pinned, not enumerated.
This one rule is what keeps \texttt{Plus[x\_\_, y\_\_]} from hanging on a wide sum. The
second is that for an ordinary (ordered) head the consumed run and the remainder are
handed to the recursion as \emph{slices} of the caller's array---no copying.

\texttt{Orderless} heads are where the real cost lives, because a commutative match may
pair a pattern with arguments in any position, not just a contiguous run. There the
matcher enumerates \(k\)-element \emph{subsets} of the arguments, which is genuinely
combinatorial; it prunes hard (moving fixed elements ahead of variable-length blanks
so the selective patterns filter first), but the worst case is why matching a
complicated pattern against a very large sum can be slow. \texttt{Flat} heads add a
third wrinkle: a single pattern element may consume several arguments, which the
matcher assembles back into an \texttt{f[\dots]} sub-expression to match against.

\subsection*{When the matcher calls the evaluator}

A \mcode{PatternTest} (\(\mathtt{p\,?\,test}\)) and a \mcode{Condition}
(\(\mathtt{p\,/;\,c}\)) are not structural: they are predicates. When the matcher reaches one, it substitutes
the current bindings into the test or condition and \emph{calls
\texttt{evaluate}}---the matcher and the evaluator are mutually recursive. A test that
does not return \texttt{True} fails the match and triggers the same
\texttt{env\_rollback} as a structural mismatch, so the partition search moves on to
the next candidate. This is why \texttt{\_?EvenQ} and \texttt{x\_ /; x > 5} work: the
condition is arbitrary Mathilda code, run mid-match with the tentative bindings in
place.

One attribute is a matcher-only concern: \texttt{OneIdentity}. It is \emph{not}
applied during evaluation (\(\mathtt{f[x]}\) is never silently collapsed to
\texttt{x}); it exists solely so that, when a structural match of \texttt{f[\dots]}
fails, the matcher may retry treating the whole subject as a single argument---the
behaviour that lets a pattern like \(\mathtt{a\_.\, x\_ + b\_.}\) match a bare
\texttt{x}. Once a match succeeds, \texttt{replace\_bindings} substitutes the captured
values into the rule's right-hand side, splicing any \texttt{Sequence} bindings into
place, and returns the rewritten expression.

\section{Rules and the symbol table}
\label{sec:int-symtab}

A rule is stored, retrieved and applied through the symbol table, and the design of
that table is what makes rule dispatch fast even for a symbol with hundreds of
definitions. Every symbol resolves to a \texttt{SymbolDef}: its \texttt{OwnValues}
(bare assignments like \texttt{x = 5}), its \texttt{DownValues} (rules on a head, like
\(\mathtt{f[x\_] := x^2}\)), an optional builtin C function pointer, its attribute
bits, its docstring, and its NDArray kernels. A rule is a \texttt{Rule} node:

\begin{ccode}[src/symtab.h (abridged)]
typedef struct Rule {
    Expr* pattern;
    Expr* replacement;
    int32_t     dispatch_arity;        /* arg count, or -1 for variable arity   */
    const char* first_arg_head_canon;  /* interned head of arg 1, or NULL        */
    int32_t     specificity;           /* larger = more specific; sorts the list */
    struct Rule* next;
} Rule;
\end{ccode}

\subsection*{Specificity, not recency}

\texttt{DownValues} are kept in a linked list, and the order of that list is the order
the matcher tries them---so it decides which of several applicable rules wins. The
ordering is by \emph{specificity}: a more specific pattern (a literal
\texttt{fib[0]}) sorts ahead of a more general one (\texttt{fib[n\_]}), with insertion
order breaking ties. It is a common misconception that new definitions simply go to
the front of the list. They do not. A general rule entered \emph{after} a specific one
still yields to it:

\mtranscript{10-internals/matcher-downvalues}

The Fibonacci base cases \texttt{fib[0]} and \texttt{fib[1]}, though entered after the
recursive rule, are tried first, and \texttt{DownValues[fib]} shows them at the front
of the list. The decisive case is \texttt{r}: the specific \texttt{r[0]} is entered
first, the general \texttt{r[x\_]} second, and yet \texttt{r[0]} still returns
\texttt{special}---specificity beat recency.

\usagebox{DownValues}

Before the matcher is ever invoked, a cheap filter discards rules that cannot possibly
apply. The \texttt{dispatch\_arity} and \texttt{first\_arg\_head\_canon} fields on each
rule are a precomputed dispatch key: a rule whose pattern takes three arguments is
skipped outright for a two-argument call, and a rule whose first argument must have
head \texttt{Plus} is skipped when the input's first argument is an
\texttt{Integer}---an interned-pointer comparison, no matching attempted. Only the
survivors reach the backtracking matcher of \S\ref{sec:int-match}.
\texttt{OwnValues} get an even shorter path: the overwhelmingly common
\texttt{x = value} binding is recognised and returned without building a match
environment at all.

\subsection*{Interning and contexts}

The symbol table is also the string interner. Every symbol name exists exactly once as
a canonical \texttt{const char*}, so two references to \texttt{Sin} share one pointer
and identity tests reduce to pointer equality---which is why the evaluator's hot paths
can write \texttt{head == SYM\_List} instead of a string comparison. Those
\texttt{SYM\_*} constants are the interned pointers for the kernel's own symbols,
filled in once at startup.

Layered over interning is the \B{Context} system: symbols live in namespaces, and
\(\mathtt{\$Context}\) is where new ones are created.

\mtranscript{10-internals/symtab-context}

A fresh symbol is born in \texttt{Global`}; a builtin lives in \texttt{System`};
\B{Begin} switches the current context so that symbols mentioned inside it are
namespaced (\texttt{here} becomes \texttt{MyPkg`here}), and \B{End} restores the
previous context. This is the machinery packages are built on.

The rule \emph{engine}---\B{Replace}, \B{ReplaceAll} (\texttt{/.}) and
\B{ReplaceRepeated} (\texttt{//.})---applies rules without installing them on a symbol:

\mtranscript{10-internals/rules-replace}

Note the semantics the internals enforce. \texttt{/.} rewrites top-down and does
\emph{not} re-descend into a subexpression it has just rewritten---so
\texttt{f[f[x]] /. f[a\_] -> a} yields \texttt{f[x]}, not \texttt{x}. \texttt{//.}
repeats to a fixed point and does reach \texttt{x}. \B{Replace} acts only at the top
level by default, which is why it leaves the list of \texttt{f}'s untouched. And
\B{ReplaceList} returns \emph{every} way the pattern can match---the direct use of the
matcher's success callback, enumerating all six orderings of the commutative sum.

\section{Memory management}
\label{sec:int-memory}

Mathilda manages memory by hand---there is no garbage collector---and the entire
scheme rests on one fact you have already met: \texttt{expr\_copy} is not a deep copy.
It is a reference-count bump.

\begin{ccode}[src/expr.c (abridged)]
Expr* expr_copy(Expr* e) {          /* a logical copy: no tree is duplicated */
    if (!e) return NULL;
    e->refcount++;
    return e;
}

void expr_free(Expr* e) {
    if (!e) return;
    if (e->refcount > 1) { e->refcount--; return; }  /* still shared: keep it */
    /* last reference: release the payload by type -- mpz_clear a bignum, free
     * an NDArray's buffer, recurse into a function node's head and args -- then
     * recycle the node struct. Interned symbol names are never freed. */
}
\end{ccode}

Sharing a subtree costs one increment; freeing decrements and reclaims only on the
last reference. The invariant that makes this safe is that a node with
\texttt{refcount > 1} is \emph{immutable}---any code wanting to rewrite a shared node
first calls \texttt{expr\_unshare}, which returns a private one-level copy (its
children still shared) and resets the metadata caches. This copy-on-write discipline
is what lets the cached hash and evaluation stamp of \S\ref{sec:int-expr} be trusted.

\subsection*{The builtin ownership contract}

Every builtin obeys one contract, and getting it right is the single most important
rule for anyone extending the system. A builtin \emph{takes ownership} of its argument
expression. It returns either a freshly built \texttt{Expr*}---in which case the
evaluator frees the original---or \texttt{NULL} to mean ``not my case,'' in which case
the evaluator keeps the expression and leaves it unevaluated for a rule to handle. That
\texttt{NULL} convention is exactly what lets a symbolic argument flow untouched
through a head that only knows how to handle numbers. A builtin that wants to reuse a
piece of its input must null the slot out before returning it, so the evaluator's
subsequent free does not double-free the reused child.\calloutnote{Under the hood}{Two
free-lists keep allocation cache-hot. Fixed-size \texttt{Expr} node structs are pooled
in a bounded list (the dead node's link is stored in a union field it no longer needs),
and the small argument arrays are pooled in per-length buckets for lengths 1 through 8
---which is why the argument array is sized exactly, never over-allocated. Both pools
are drained by an \texttt{atexit} hook, so a clean run leaves no reachable-but-unfreed
memory and \texttt{valgrind} stays quiet.}

You can watch the footprint from the language. \B{MemoryInUse} and \B{MaxMemoryUsed}
report the process heap in bytes:

\mtranscript{10-internals/memory-inuse}

Building a million-element \B{Range} raises the footprint by about eight
megabytes---a dense \texttt{int64} buffer of \(10^6\) elements
(\S\ref{sec:int-packed}), not a million boxed nodes. \B{Clear} releases the binding,
but the resident figure does not fall: freed memory returns to the process's own pools
and allocator, not immediately to the operating system, so \B{MaxMemoryUsed} records
the high-water mark. The numbers here are re-measured on every build and depend on the
machine; it is their shape that is the lesson.

\section{The parser and the printer}
\label{sec:int-parse}

The parser and printer are the two ends of the pipeline---the translation between the
string a human types and the tree the evaluator rewrites.

The parser (\texttt{src/parse.c}) is a \emph{Pratt}, or precedence-climbing, parser
with no separate lexer: tokens are read inline as parsing proceeds. Every operator
sits on a numeric precedence ladder that runs from 1 to 10000 in round steps, with
generous gaps so a new operator can be slotted between two existing ones without
renumbering. The best way to see the parser's decisions is to X-ray them with
\B{FullForm}, which shows the tree the grouping produced:

\mtranscript{10-internals/parser-fullform}

Every line is a precedence decision made visible. \texttt{a + b*c} groups the
multiplication first; \texttt{a\^{}b\^{}c} associates to the right; subtraction and
division are not primitive at all---\texttt{a - b} is \texttt{Plus[a, Times[-1, b]]}
and \texttt{a/b} is \texttt{Times[a, Power[b, -1]]}, so the evaluator needs rules only
for \texttt{Plus}, \texttt{Times} and \texttt{Power}. A chain of equalities folds into
a single \texttt{Inequality} node; \texttt{\&\&} binds tighter than \texttt{||}; the
prefix \texttt{\#\^{}2 \&} is \texttt{Function[Power[Slot[1], 2]]}; and
\texttt{a[[1]]} is \texttt{Part[a, 1]}. Implicit multiplication---the space in
\texttt{b c}---is injected by the parser as an ordinary \texttt{Times} at the
multiplication rung.

The printer (\texttt{src/print.c}) runs the same ladder in reverse, adding parentheses
exactly where a subexpression's precedence is lower than its context so that the output
re-parses to the same tree:

\mtranscript{10-internals/printer-forms}

\B{InputForm} prints a form suitable for re-reading, \B{FullForm} the raw tree, and
the default \texttt{OutputForm} the familiar two-dimensional-ish notation. The printer
restores \texttt{Times[a, Power[b, -1]]} to \texttt{a/b}, \texttt{Rational[3, 4]} to
\texttt{3/4}, and parenthesises \texttt{(a + b) c} because \texttt{Plus} binds looser
than \texttt{Times}. It is the same precedence table the parser used, which is what
guarantees the round trip.

\section{The read--eval--print loop and session state}
\label{sec:int-repl}

The evaluator is wrapped by the interactive loop (\texttt{src/repl.c}), and it is
worth being precise that the loop is a distinct layer---it does several things
\emph{around} each call to \texttt{evaluate} that the evaluator itself knows nothing
about. For one input line, in order, the loop:

\begin{enumerate}
  \item sets \(\mathtt{\$Line}\) to the current line number;
  \item applies the \(\mathtt{\$PreRead}\) hook to the raw input \emph{text};
  \item parses the text to an \texttt{Expr};
  \item stores it as \texttt{In[\textit{n}]} (a \texttt{DownValue} on \texttt{In}), so
        the literal input can be recalled;
  \item applies the \(\mathtt{\$Pre}\) hook to the parsed expression;
  \item calls \texttt{evaluate};
  \item applies the \(\mathtt{\$Post}\) hook to the result;
  \item stores the result as \texttt{Out[\textit{n}]} (a \texttt{DownValue} on
        \texttt{Out});
  \item applies the \(\mathtt{\$PrePrint}\) hook to a copy for display;
  \item prints \texttt{Out[\textit{n}]= } and the expression.
\end{enumerate}

The history reach-backs are built on step 8: \texttt{\%} parses to
\texttt{Out[-1]}, \texttt{\%\%} to \texttt{Out[-2]}, and \texttt{Out[\textit{n}]}
resolves through the same \texttt{DownValue} store, a negative index counting back from
\(\mathtt{\$Line}\). Because all of this lives in the loop and not in \texttt{evaluate},
it is present only in an interactive session. The proof is to run the same lines
\emph{non-interactively}---which is exactly how every transcript in this book is
produced, through a batch pipe:

\mtranscript{10-internals/repl-history}

Outside the interactive loop there is no \texttt{In}/\texttt{Out} store and no running
\(\mathtt{\$Line}\): \texttt{\%} stays the symbolic \texttt{Out[-1]}, \texttt{Out[1]}
does not resolve, and \(\mathtt{\$Line}\) is just an unbound symbol. Nothing is broken
---this is the same evaluator---but the session bookkeeping simply is not there,
because the loop that maintains it is not running. That is the clearest possible
demonstration that history is a service of the REPL, layered on top of a stateless
evaluator.\calloutnote{Under the hood}{The loop chooses one of three modes at startup.
A script given with \texttt{-file} is evaluated statement by statement with nothing
echoed. A non-terminal standard input (a pipe) selects a line-oriented NDJSON protocol
---the one that drives this book's build. Only a real terminal runs the interactive
loop, with GNU Readline and the completeness-driven multi-line editing described in
\texttt{SPEC.md}~\S11, where pressing Return submits only when brackets, strings and
comments are all balanced.}

\section{Numbers and the packed-array substrate}
\label{sec:int-packed}

Numbers deserve a closer look, because ``number'' is really several internal
representations that the system chooses between automatically. The
\texttt{Head}, though, is uniform:

\mtranscript{10-internals/numbers-tower}

An exact integer is a machine \texttt{int64} promoted to a GMP bignum as needed
(\S\ref{sec:int-expr}); a rational is a \texttt{Rational} of two integers; a machine
real is a \texttt{double}; and at higher precision the arithmetic moves onto MPFR,
which is why \texttt{Precision[N[Pi, 40]]} reports about forty digits rather than the
\texttt{MachinePrecision} of a bare \texttt{N}. Exactness is contagious in the usual
way: \texttt{1/2 + 1/3} stays the exact \texttt{5/6}, but a single machine real, as in
\texttt{1/2 + 0.5}, makes the whole result inexact. This numeric tower is the subject
of \S\ref{sec:numcalc} from the mathematics side; here it is enough to see that one
\texttt{Head} hides several machine types.

The last representation is the one that makes numeric work fast: the packed array. A
list all of whose elements are machine numbers can be stored not as a tree of boxed
nodes but as a single \texttt{EXPR\_NDARRAY}---a dense, flat buffer. Crucially this is
\emph{one node type with two surfaces}. A visible \texttt{NDArray[\dots]} that a user
typed is atomic: its head is \texttt{NDArray}, and \texttt{AtomQ} is true. A
\emph{packed list} is the same buffer wearing the head \texttt{List}: it is
indistinguishable from the nested list it stores---\texttt{ListQ} is true, patterns
match it, it prints as a list---and only \B{NDArrayQ} can tell the two apart.

\mtranscript{10-internals/packed-gate}

That indistinguishability is a promise the evaluator has to keep, and it keeps it with
the \emph{transparency gate} of \S\ref{sec:int-eval}, step~7.5. There are thousands of
places in the source that read a function node's argument array directly; a packed
buffer has no such array. So before the evaluator hands a packed argument to any head
that has not explicitly declared itself packed-aware, it \emph{materializes} the buffer
back into an ordinary list. A head opts in only after it has been checked to read the
buffer correctly, and the gate additionally materializes an integer buffer for any head
not proven exact on \texttt{int64}---which is why \texttt{Total[Range[5]]} stays the
exact integer \texttt{15} and never becomes \texttt{15.}. The full story of packing,
the exactness contract, and the audit that guards the list of aware heads is in
Chapter~\ref{ch:datastructures} (\S\ref{sec:ds-packed}); what matters here is that
packing is a performance representation the evaluator makes invisible, at the cost of
one gate check per call.

\section{Bootstrapping: startup and the two-tier design}
\label{sec:int-boot}

When Mathilda starts, \texttt{main()} in \texttt{src/repl.c} builds the world in a
fixed order: \texttt{symtab\_init()} creates the symbol table and interns the kernel's
\texttt{SYM\_*} names, then \texttt{core\_init()} in \texttt{src/core.c} registers the
core builtins and calls each subsystem's \texttt{*\_init()} function in turn---several
dozen of them, from arithmetic through the special functions and the packed-array
substrate to the higher-level algebra and graphics services. That ordered list of calls
in \texttt{core.c} is the authoritative source of truth for initialization order.

Then comes the seam that gives the system its character. After the C side is built,
\texttt{main()} loads \texttt{src/internal/init.m}, which in turn reads the remaining
bootstrap \texttt{.m} files. This is Mathilda's \emph{two-tier design}:
performance-sensitive primitives are written in C, while a great deal of the actual
mathematics is expressed as ordinary rewrite rules in Mathilda's own language. The
derivative rules live in \texttt{deriv.m}; the integral tables are \texttt{.m} data;
and the same fixed-point evaluator that runs your rules runs these:

\mtranscript{10-internals/twotier-rules}

\texttt{D[Sin[x], x]}, the chain rule for \texttt{D[f[g[x]], x]}, and the
antiderivative of \(1/(1+x^2)\) are not special-cased in C---they are
\texttt{DownValues} loaded at startup, tried at step~11 of every relevant evaluation
(\S\ref{sec:int-eval}) exactly as a rule you write would be. This is the deepest
reason the architecture holds together: because everything is an expression and the
evaluator is a generic fixed-point loop, extending the system's mathematics most often
means adding a rule, not writing C. The recipes for doing so---adding a builtin, a
module, a pattern, an operator---are Chapter~\ref{ch:contributing}, and the build,
test and audit machinery that keeps all of this honest is
Chapter~\ref{ch:development}.
